\ficha{Hanke e Schuler (2015), Currency Boards for Developing Countries}{hankeschuler2015}

\begin{identificacao}
HANKE, Steve H.; SCHULER, Kurt. \emph{Currency Boards for Developing
Countries: A Handbook}. Edição revista. [S.l.]: os autores, 2015. 86 p. A edição de 2015 não
traz local nem editora: é reedição autorizada, distribuída pelos autores.
Edição original: San Francisco: ICS Press, para o International Center for
Economic Growth, 1994.

Manual normativo de economia monetária aplicada, em inglês, lido em leitura
seletiva: capítulos 1, 2, 3, 6 e 7, cerca de cinquenta das oitenta páginas
numeradas. Os capítulos 4 e 5, que são o passo a passo de implantação e
operação, foram pulados. A versão lida é a revista de 2015, com prefácio
datado de julho de 2015, que atualiza referências a acontecimentos e,
segundo os próprios autores, não é revisão integral do texto de 1994.
\end{identificacao}

\bloco{Problema e tese}
Como um país em desenvolvimento obtém moeda sã, isto é, estável, crível e
plenamente conversível (p. 1)? A resposta dos autores é que o banco central
típico não a produz e a caixa de conversão produz. O argumento não é que a
regra funcione melhor em média, e sim que a discricionariedade monetária, em
país de finanças públicas frouxas, sempre acaba usada para financiar déficit,
de modo que a única forma de tornar a taxa fixa crível é retirar da autoridade
emissora a capacidade técnica de descumpri-la.

\bloco{Argumento}
O livro inteiro é um manual de reforma. Ele defende a caixa de conversão
(capítulos 1 e 2), expõe a mecânica de oferta de moeda e a história do
instrumento (capítulo 3), ensina passo a passo como converter um banco central
em caixa ou como criar uma moeda paralela (capítulos 4 e 5), responde a oito
objeções (capítulo 6) e fecha com uma constituição-modelo de caixa de conversão
em apêndice. É literatura de advocacia técnica escrita por dois assessores de
reformas monetárias, não historiografia.

\begin{enumerate}[leftmargin=*,nosep]
\item A definição é institucional e curta.
\chave{A currency board is a monetary institution that issues notes and coins
... fully backed by a foreign `reserve' currency and fully convertible into the
reserve currency at a fixed rate and on demand}{p. 2}
O banco central, por contraste, é definido como autoridade monetária com
controle discricionário e monopolista da oferta de reservas dos bancos
comerciais (p. 2).
\item A tabela comparativa (Figura 1.1, p. 3) alinha dezesseis pares de
características. As decisivas: lastro externo de 100 por cento contra reservas
variáveis, taxa fixa contra taxa atrelada ou flutuante, política sujeita a
regra contra política discricionária, ausência de emprestador de última
instância, e impossibilidade de financiar o governo. A caixa é passiva: o
mercado determina a quantidade emitida (p. 5).
\item O lastro é de 100 por cento ou pouco mais, com margem histórica de 105 a
110 por cento como colchão contra perda de valor dos títulos (p. 4). Só as
obrigações da caixa são lastreadas assim; os bancos comerciais não precisam
guardar 100 por cento, e apenas a base monetária, não M1 ou M2, tem cobertura
integral (p. 4-5).
\item A distinção entre taxa fixa e taxa atrelada é o núcleo analítico. Atrelada
é a taxa constante por ora, sem garantia crível de permanência; fixa é a
permanente, alterável apenas por regra conhecida de antemão (p. 15). Os autores
registram que o único período em que muitos bancos centrais mantiveram taxas
verdadeiramente fixas foi o padrão-ouro clássico, de 1880 ou antes até 1914,
quando os bancos centrais eram minoria e muitos eram privados (p. 15).
\item Há um teste de diagnóstico explícito, no trecho sobre bancos centrais
confundidos com caixas de conversão. O que importa não é a razão entre reservas
externas e notas em circulação, e sim entre reservas externas e o total das
obrigações do emissor. Uma autoridade com 100 por cento contra notas mas não
contra depósitos não é caixa de conversão (p. 38-39).
\item A história é apresentada como síntese de estudos anteriores. Mais de
setenta países tiveram caixas de conversão; a primeira foi a das Maurícias, em
1849.
\chave{Currency boards spread slowly until about 1900, when a few other British
colonies and Argentina, an independent country, established currency
boards}{p. 40}
O auge foi por volta de 1950, na África, no Caribe e no sul da Ásia. Todas
mantiveram taxa fixa e conversibilidade plena, inclusive durante a Grande
Depressão (p. 40-41). A conversão em bancos centrais, no fim dos anos 1950 e
nos anos 1960, é atribuída a política, não a desempenho.
\chave{A central bank was a symbol of independence, like a national flag}{p. 41}
\item Os autores separam caixas ortodoxas dos sistemas ``currency board-like''
dos anos 1990, que diferem em dois pontos: não têm teto de reserva e podem
atuar como emprestador de última instância (p. 42-43). A Argentina de 1991 a
2002 entra na tabela de casos, mas em nota os autores dizem incluí-la apenas
porque se supõe comumente que ela tenha tido uma caixa de conversão, e remetem
a seus próprios textos sobre por que ela não teve (p. 43).
\end{enumerate}

\bloco{Conceitos e definições operacionais}
Moeda sã tem três atributos operacionais: estabilidade, ou inflação anual
corrente de um dígito; credibilidade, ou confiança de que a inflação futura
seguirá baixa; conversibilidade plena, ou compra irrestrita de bens e de
moedas estrangeiras a preço de mercado (p. 1). Moeda de reserva é a moeda
estrangeira ou mercadoria escolhida pela estabilidade esperada, e o país que a
emite é o país de reserva (p. 2). Lastro integral é ativo pagável em moeda de
reserva equivalente a 100 por cento das obrigações da caixa, fixado em lei
(p. 4). Senhoriagem, para a caixa, é só o juro dos títulos de reserva, nunca a
inflação (p. 7). Restrição orçamentária dura é a impossibilidade de gastar
acima do que se arrecada e se toma emprestado sem subsídio (p. 15). Ortodoxia é
o conjunto completo das características da Figura 1.1, e a falta de qualquer
uma delas define o sistema apenas semelhante a uma caixa de conversão
(p. 42-43).

\bloco{Evidência e método}
Não há teste estatístico, regressão nem contrafactual. O capítulo histórico é
declaradamente síntese de estudos anteriores sobre as caixas coloniais
britânicas, listados em nota (p. 40), e a comparação entre caixas e os bancos
centrais que as sucederam vem de uma única fonte quantitativa, Schuler (1992b:
202-3), citada para inflação mais alta e crescimento mais baixo depois da
conversão (p. 42). A Figura 3.13 (p. 43) tabula os casos vigentes e recentes,
com população, produto, datas e taxa. A comparação é entre instituições
``típicas'', e os autores avisam que não descrevem tipos ideais (p. 3). Eles
admitem que não existe estudo sistemático sobre choques de câmbio real em
caixas de conversão (p. 41).

\bloco{Agentes e vertentes}
Hanke e Schuler são partidários e praticantes: Hanke assessorou Argentina,
Estônia, Bulgária, Bósnia, Equador, Lituânia e Montenegro (p. vii). A filiação
é ao liberalismo monetário do Cato Institute, com simpatia declarada pelo
\emph{free banking} de Dowd, Selgin e White (p. 9). Os adversários nomeados são
o Fundo Monetário Internacional, por exigir precondições fiscais para a reforma
(Baliño, Enoch e Gulde, 1997, p. 7), Williamson (1995) e Roubini (1998), por
não distinguirem caixas ortodoxas de sistemas apenas semelhantes (p. 45).
Friedman aparece como defensor do câmbio flutuante para países desenvolvidos e
da taxa fixa para os demais (p. 67). A literatura colonial britânica, Chalmers,
Clauson, Sayers e King, é a base do capítulo histórico.

\bloco{Críticas e limites}
O livro é normativo e assume o que deveria demonstrar. O desempenho histórico
das caixas é medido contra bancos centrais posteriores dos mesmos países, sem
controle para época, choque externo ou administração colonial. Os próprios
autores registram que as caixas coloniais não tinham proteção legal formal e
sim proteção informal, por serem geridas por funcionários britânicos e sob
ameaça de demissão por Londres (p. 40), o que torna impossível separar o
efeito da instituição do efeito do governo colonial. O critério de ortodoxia
funciona ex post como blindagem: o caso que fracassou, a Argentina, é retirado
da amostra por não ter sido ortodoxo (p. 43). Sobre o Brasil o livro nada diz.
O país aparece uma única vez, como maior parceiro comercial da Argentina na
crise de 1998 (p. 44). A Caixa de Conversão não é mencionada, e a Caja de
Conversión argentina anterior a 1914 não é sequer nomeada: o caso argentino do
começo do século aparece em duas orações (p. 40 e p. 70).

\begin{dialogo}
\item Procurar nos jornais se o desenho da Caixa foi discutido em termos de
lastro integral e de proibição de emprestar ao Tesouro ou ao Banco do Brasil.
Testa se o debate brasileiro produziu por conta própria os requisitos que
Hanke e Schuler tratam como definidores (p. 2 e p. 7), ou se tratou a Caixa
como repartição do Tesouro.
\item Procurar o vocabulário sobre limite de emissão e fundo de garantia, e
quem propõe baixar a cobertura. Os autores sustentam que razão inferior a 100
por cento gera pressão política para reduzi-la outra vez (p. 70). Testa se essa
dinâmica aparece na imprensa entre 1906 e 1913, quando o limite foi ampliado.
\item Procurar quem atribui à Caixa funções de socorro bancário ou de
redesconto. Testa se a imprensa separava a Caixa do banco emissor, fronteira
que os autores tratam como a que distingue caixa de conversão de banco central
(p. 5 e p. 64-65).
\item Procurar se a taxa de 15 pence é defendida como irrevogável ou como
provisória. Testa a distinção entre taxa fixa e taxa atrelada (p. 15) e permite
datar quando o compromisso brasileiro passou a ser lido como revogável, o que
de fato ocorreu em 1910 e em 1914.
\item Procurar menções à experiência argentina da Caja de Conversión. Hanke e
Schuler registram a Argentina como caso independente do começo do século
(p. 40 e p. 70) mas não a descrevem. Se os jornais brasileiros discutiam o caso
argentino, a comparação da Parte IV sai das fontes, e não desta literatura.
\end{dialogo}

\usonocapitulo{Parte IV, e só para explicitar a definição moderna de caixa de
conversão e sua história institucional. O texto não descreve a Caixa de
Conversão brasileira e não autoriza presumir que ela obedecesse a esse modelo:
o livro não trata do Brasil, não menciona a Caixa e cita o país uma única vez,
como parceiro comercial da Argentina em 1998 (p. 44). O que ele fornece é
critério, não veredito, e o critério é duplo: a lista de características da
Figura 1.1 (p. 3) e o teste da razão entre reservas externas e o total das
obrigações do emissor, e não entre reservas e notas em circulação (p. 38-39).
Aplicar esse teste ao arranjo de 1906, com as fontes e a legislação, é trabalho
da dissertação, e o resultado pode ser negativo. A categoria intermediária dos
autores, sistema apenas semelhante a uma caixa de conversão (p. 42-43), serve
para classificar sem forçar. Fica também o registro de que a definição atual do
instrumento é retrospectiva, formulada nos anos 1990 sobre a experiência
colonial britânica, e que os agentes de 1906 não dispunham dela.}
